\abstracttext{As audiências públicas da Câmara dos Deputados chegam ao público sobretudo por matérias
jornalísticas que selecionam poucas opiniões, sem indicar em que ponto da transcrição cada uma foi
dita. Inspirado nas trilhas de auditoria documental e de mineração de opiniões, este trabalho propõe
um pipeline sobre o PublicHearingBR com três componentes. A Unidade Deliberativa Verificável liga cada
opinião publicada a uma sentença da fala da própria pessoa, com a posição na transcrição, o tipo e a
origem do vínculo; das 2.203 opiniões, 2.105 recebem evidência. Os perfis por ator são escritos por
um modelo de linguagem somente a partir das falas de cada pessoa em audiências de treino, em blocos
de posições, critérios, alinhamentos declarados e forma de argumentar. Na simulação, o perfil orienta
um modelo de linguagem a responder como a pessoa responderia, com trechos recuperados de suas falas
anteriores e uma técnica de decodificação que amplifica a influência desse material
(\textit{classifier-free guidance}); cada resposta declara a base da estimativa e traz uma
justificação cujas cópias literais são conferidas automaticamente. A avaliação usa as unidades
verificáveis de audiências posteriores como gabarito de perguntas de múltipla escolha com distratores
da mesma audiência, lidas pelas probabilidades das letras em quatro ordens das opções, com
hiperparâmetros escolhidos na validação e intervalos obtidos por reamostragem de audiências.
<Resultados principais: acerto de cada condição no teste e diferenças entre condições com intervalo
de confiança.>}
